% !TeX program = pdflatex
\documentclass[aspectratio=169,xcolor={dvipsnames,table}]{beamer}

\input{../../_en_preambulo_beamer}
% Comandos y macros estadísticas compartidas para las presentaciones Beamer en inglés (Metropolis)

% Conjuntos numéricos básicos
\newcommand{\R}{\mathbb{R}}
\newcommand{\N}{\mathbb{N}}
\newcommand{\Q}{\mathbb{Q}}
\newcommand{\Z}{\mathbb{Z}}
\newcommand{\set}[1]{\left\{ #1 \right\}}
\newcommand{\abs}[1]{\left|#1\right|}
\newcommand{\norm}[1]{\left\| #1 \right\|}

% Comandos estadísticos y probabilísticos del curso
\newcommand{\Var}{\operatorname{Var}}
\newcommand{\cov}{\operatorname{Cov}}
\newcommand{\corr}{\rho}
\newcommand{\comb}[2]{\begin{pmatrix} #1 \\ #2 \end{pmatrix}}
\newcommand{\s}{\sigma}
\newcommand{\particion}{\mathfrak{P}}
\newcommand{\card}[1]{\#\left( #1 \right)}
\newcommand{\seq}[3]{\set{#1}_{#2}^{#3}}
\newcommand{\E}{\mathbb{E}}
\newcommand{\Prob}{\mathbb{P}}

% Entornos pedagógicos y cajas en inglés académico natural (soportando tanto nombres en inglés como en español)
\newenvironment{discussion}{%
  \begin{alertblock}{Discussion Question}%
}{%
  \end{alertblock}%
}
\newenvironment{discusion}{%
  \begin{alertblock}{Discussion Question}%
}{%
  \end{alertblock}%
}

\newenvironment{reflection}{%
  \begin{alertblock}{Classroom Reflection}%
}{%
  \end{alertblock}%
}
\newenvironment{reflexion}{%
  \begin{alertblock}{Classroom Reflection}%
}{%
  \end{alertblock}%
}

\newenvironment{paradox}{%
  \begin{alertblock}{Exploratory Paradox}%
}{%
  \end{alertblock}%
}
\newenvironment{paradoja}{%
  \begin{alertblock}{Exploratory Paradox}%
}{%
  \end{alertblock}%
}

\newenvironment{motivation}{%
  \begin{exampleblock}{Motivation: Data Science in Practice}%
}{%
  \end{exampleblock}%
}
\newenvironment{motivacion}{%
  \begin{exampleblock}{Motivation: Data Science in Practice}%
}{%
  \end{exampleblock}%
}

\newenvironment{quickchallenge}{%
  \begin{block}{Quick Team Challenge}%
}{%
  \end{block}%
}
\newenvironment{retoclase}{%
  \begin{block}{Quick Team Challenge}%
}{%
  \end{block}%
}


\title{Point Estimation, Unbiasedness, Efficiency, and Consistency}
\subtitle{Section 06.01 --- Bias, Mean Squared Error, and the Cramer-Rao Bound}
\author[J. Castillo Colmenares]{Juliho Castillo Colmenares}
\institute[Tec de Monterrey]{Tecnologico de Monterrey}
\date{\vspace{-1.5cm}}

\begin{document}

% SLIDE 01: Title Page
\begin{frame}[plain]
	\titlepage
\end{frame}

% SLIDE 02: Roadmap
\begin{frame}{Roadmap --- Chapter 06: Estimation and Data Science}
	\vspace{-0.15cm}
	\begin{block}{Curricular Structure of Unit 5}
		\footnotesize
		\begin{enumerate}\itemsep=0.03cm
			\item \textbf{\textcolor{TecRojo}{Section 06.01: Point Estimation, Unbiasedness, Efficiency, and Consistency}} $\leftarrow$ \emph{Current focus}
			\item \textcolor{gray}{Section 06.02: Method of Moments (MoM)}
			\item \textcolor{gray}{Section 06.03: Maximum Likelihood Estimation (MLE) and Score}
			\item \textcolor{gray}{Section 06.04: Confidence Intervals for Population Means ($Z$ and $t$)}
			\item \textcolor{gray}{Section 06.05: Confidence Intervals for Variances ($\chi^2$) and Proportions}
		\end{enumerate}
	\end{block}
	\vspace{-0.2cm}
	\begin{alertblock}{Learning Objective}
		\scriptsize
		Formalize the criteria that define a ``good'' estimator --- bias, mean squared error, relative efficiency, and consistency --- and establish the Cramer-Rao bound as the theoretical precision limit no unbiased estimator can beat.
	\end{alertblock}
\end{frame}

% SLIDE 03: Motivation
\begin{frame}{From Sampling Distributions to Estimation Theory}
	\footnotesize
	\begin{motivacion}
		In Chapter 05 we studied the sampling distribution of $\bar X$ and $S^2$. Now we take a step back: what makes an estimator ``good''? Before building estimators (Sections 06.02--06.03), we need a formal framework for \emph{evaluating} them.
	\end{motivacion}
	\vspace{0.15cm}
	\pause
	\begin{itemize}\itemsep=0.1cm
		\item \textbf{Unbiasedness:} does the estimator hit the target ``on average''? \pause
		\item \textbf{Efficiency:} among several unbiased estimators, which has lower variance? \pause
		\item \textbf{Consistency:} does the estimator improve as the sample grows?
	\end{itemize}
\end{frame}

% SLIDE 03B: Bias and MSE
\begin{frame}{Bias and Mean Squared Error}
	\vspace{-0.15cm}
	\begin{columns}[T]
		\begin{column}{0.48\textwidth}
			\begin{block}{Definitions}
				\footnotesize
				$\text{Bias}(\hat\theta) = E(\hat\theta)-\theta$, $\text{MSE}(\hat\theta) = E[(\hat\theta-\theta)^2]$.
			\end{block}
		\end{column}
		\begin{column}{0.48\textwidth}
			\begin{alertblock}{Bias-Variance Decomposition}
				\footnotesize
				$\text{MSE}(\hat\theta) = \Var(\hat\theta) + [\text{Bias}(\hat\theta)]^2$. A slightly biased estimator can have lower MSE than an unbiased one.
			\end{alertblock}
		\end{column}
	\end{columns}
\end{frame}

% SLIDE 04: Efficiency and Cramer-Rao
\begin{frame}{Relative Efficiency and the Cramer-Rao Bound}
	\vspace{-0.15cm}
	\begin{columns}[T]
		\begin{column}{0.48\textwidth}
			\begin{block}{Relative Efficiency}
				\footnotesize
				$\text{Eff}(\hat\theta_1,\hat\theta_2) = \Var(\hat\theta_2)/\Var(\hat\theta_1)$ for unbiased estimators.
			\end{block}
		\end{column}
		\begin{column}{0.48\textwidth}
			\begin{alertblock}{Cramer-Rao Lower Bound (CRLB)}
				\footnotesize
				$\Var(\hat\theta) \geq 1/(nI(\theta))$. An estimator attaining this bound is a UMVUE: the best possible.
			\end{alertblock}
		\end{column}
	\end{columns}
\end{frame}

% SLIDE 05: Consistency
\begin{frame}{Consistency}
	\vspace{-0.15cm}
	\begin{block}{Definition}
		$$\hat\theta_n \xrightarrow{P} \theta \iff \forall\varepsilon>0,\ \lim_{n\to\infty} P(|\hat\theta_n-\theta|\ge\varepsilon) = 0$$
	\end{block}
	\pause
	\vspace{0.1cm}
	\begin{alertblock}{Practical Sufficient Condition}
		By Chebyshev, $P(|\hat\theta_n-\theta|\ge\varepsilon) \le \text{MSE}(\hat\theta_n)/\varepsilon^2$. If $\text{MSE}(\hat\theta_n)\to 0$, then $\hat\theta_n$ is consistent --- the most common way to prove it in practice.
	\end{alertblock}
\end{frame}

% SLIDE 06: Applications
\begin{frame}{Applications in Data Science}
	\vspace{-0.1cm}
	\begin{itemize}\itemsep=0.1cm
		\item \textbf{Regularization (Ridge/Lasso):} accepts controlled bias in exchange for lower variance and better generalization error. \pause
		\item \textbf{Cross-validation:} estimates a model's out-of-sample MSE. \pause
		\item \textbf{Estimator selection:} relative efficiency guides the choice between methods (MoM vs. MLE, Sections 06.02--06.03). \pause
		\item \textbf{Big data:} consistency guarantees that more data always helps, even with approximate models.
	\end{itemize}
\end{frame}

% SLIDE 07: Python Lab Block 1
\begin{frame}[fragile]{Python Lab: Bias, MSE, and Efficiency}
	\vspace{-0.05cm}
	\scriptsize
	Bias-variance decomposition and relative efficiency comparing $\bar X$ against a naive estimator (\texttt{06.01\_point\_estimation\_quality.py}):
	{\lstset{basicstyle=\fontsize{4pt}{4.8pt}\selectfont\ttfamily,aboveskip=0.1em,belowskip=0.1em}
	\lstinputlisting[firstline=17, lastline=36]{../../code/06_estimacion_estadistica/06.01_point_estimation_quality.py}}
\end{frame}

% SLIDE 08: Python Lab Block 2
\begin{frame}[fragile]{Python Lab: Shrinkage Estimator}
	\vspace{-0.05cm}
	\scriptsize
	Computing the optimal shrinkage factor $c^*$ that minimizes MSE, verified via grid search:
	{\lstset{basicstyle=\fontsize{4pt}{4.8pt}\selectfont\ttfamily,aboveskip=0.1em,belowskip=0.1em}
	\lstinputlisting[firstline=39, lastline=59]{../../code/06_estimacion_estadistica/06.01_point_estimation_quality.py}}
\end{frame}

% SLIDE 09: Python Lab Block 3
\begin{frame}[fragile]{Python Lab: Consistency via Chebyshev}
	\vspace{-0.05cm}
	\scriptsize
	Chebyshev bound vs. Monte Carlo empirical probability for the sample proportion:
	{\lstset{basicstyle=\fontsize{4pt}{4.8pt}\selectfont\ttfamily,aboveskip=0.1em,belowskip=0.1em}
	\lstinputlisting[firstline=62, lastline=82]{../../code/06_estimacion_estadistica/06.01_point_estimation_quality.py}}
\end{frame}

% SLIDE 10: Python Lab Terminal Output
\begin{frame}[fragile]{Python Lab: Terminal Output}
	\vspace{-0.1cm}
	\begin{block}{}
		\fontsize{4pt}{5pt}\selectfont
		\begin{verbatim}
=== Block 1: Bias-Variance-MSE Decomposition and Relative Efficiency ===
Xbar: bias=0.009, Var=6.24, MSE=6.24 | T1: bias=0.023, Var=100.22, MSE=100.22
Relative efficiency Ef(Xbar, T1) = 16.07 (theoretical: n=16)

=== Block 2: Shrinkage Estimator ===
c*=0.6429 | MSE(c=1)=5.00 | MSE(c*)=3.21 | reduction: 35.7%

=== Block 3: Consistency via Chebyshev ===
n=100: bound=1.00, empirical=0.319 | n=1000: bound=0.10, empirical=0.001
Minimum n for bound<=0.10: n>=1000
		\end{verbatim}
	\end{block}
\end{frame}

% SLIDE 11: Exercise Level 1 (Statement)
\begin{frame}{In-Class Exercise --- Level 1: Fundamental (1/2)}
	\vspace{-0.1cm}
	\begin{block}{Problem 6.1.1 --- Unbiasedness of Linear Combinations (Statement)}
		Let $X_1, X_2$ be iid with $E(X_i)=\mu$. Consider $\hat\theta_w = 0.3X_1+0.7X_2$. Prove that $\hat\theta_w$ is unbiased for $\mu$, and generalize the condition on the weights.
	\end{block}
	\vspace{0.15cm}
	\pause
	\begin{alertblock}{Model Setup}
		\begin{itemize}\itemsep=0.04cm
			\item Use linearity of expectation: $E(\hat\theta_w)=0.3E(X_1)+0.7E(X_2)$.
		\end{itemize}
	\end{alertblock}
\end{frame}

% SLIDE 12: Exercise Level 1 (Solution)
\begin{frame}{In-Class Exercise --- Level 1: Fundamental (2/2)}
	\vspace{-0.05cm}
	\scriptsize
	Apply linearity: \pause
	\begin{block}{Calculation}
		$$E(\hat\theta_w) = 0.3\mu + 0.7\mu = (0.3+0.7)\mu = \mu$$
	\end{block}
	\vspace{0.05cm}
	\pause
	\begin{alertblock}{Generalization}
		\scriptsize
		For $\hat\theta=\sum_i w_i X_i$ with $E(X_i)=\mu$ for all $i$: $E(\hat\theta)=\left(\sum_i w_i\right)\mu$. The estimator is unbiased if and only if $\sum_i w_i = 1$ --- the sample mean ($w_i=1/n$) is just one particular case of this family.
	\end{alertblock}
\end{frame}

% SLIDE 13: Exercise Level 2 (Statement)
\begin{frame}{In-Class Exercise --- Level 2: Operational (1/2)}
	\vspace{-0.1cm}
	\begin{block}{Problem 6.1.6 --- Consistency via Chebyshev (Statement)}
		An unbiased estimator $\hat\theta_n$ has $\Var(\hat\theta_n)=16/n$. Use Chebyshev to bound $P(|\hat\theta_n-\theta|\ge 1)$ for $n=64$, and explain why this confirms consistency.
	\end{block}
	\vspace{0.15cm}
	\pause
	\begin{alertblock}{Strategy}
		\scriptsize
		$P(|\hat\theta_n-\theta|\ge\varepsilon) \le \Var(\hat\theta_n)/\varepsilon^2$.
	\end{alertblock}
\end{frame}

% SLIDE 14: Exercise Level 2 (Solution)
\begin{frame}{In-Class Exercise --- Level 2: Operational (2/2)}
	\vspace{-0.05cm}
	\scriptsize
	Substitute into the Chebyshev bound: \pause
	\begin{block}{Calculation}
		$$P(|\hat\theta_{64}-\theta|\ge 1) \le \frac{16/64}{1^2} = 0.25$$
	\end{block}
	\vspace{0.05cm}
	\pause
	\begin{alertblock}{Interpretation}
		\scriptsize
		Since $16/n \to 0$ as $n\to\infty$ for any fixed $\varepsilon$, the Chebyshev bound tends to $0$, confirming that $\hat\theta_n \xrightarrow{P} \theta$: the estimator is consistent.
	\end{alertblock}
\end{frame}

% SLIDE 15: Exercise Level 3 (Statement)
\begin{frame}{In-Class Exercise --- Level 3: Analytical (1/2)}
	\vspace{-0.1cm}
	\begin{block}{Problem 6.1.7 --- CRLB for the Normal Mean (Statement)}
		Let $X_1,\ldots,X_n \sim N(\mu,\sigma^2)$ with known $\sigma^2$. Compute $I(\mu)$, derive the CRLB for $\Var(\hat\mu)$, and show that $\bar X$ attains it exactly.
	\end{block}
	\vspace{0.1cm}
	\pause
	\begin{alertblock}{Strategy}
		\scriptsize
		Differentiate $\ln f(x\mid\mu)$ twice with respect to $\mu$ to obtain $I(\mu)$.
	\end{alertblock}
\end{frame}

% SLIDE 16: Exercise Level 3 (Solution)
\begin{frame}{In-Class Exercise --- Level 3: Analytical (2/2)}
	\vspace{-0.05cm}
	\scriptsize
	Compute the Fisher information: \pause
	\begin{block}{Derivation}
		\scriptsize
		\begin{align*}
			\frac{\partial}{\partial\mu}\ln f = \frac{x-\mu}{\sigma^2}, \qquad \frac{\partial^2}{\partial\mu^2}\ln f = -\frac{1}{\sigma^2} \implies I(\mu) = \frac{1}{\sigma^2}.
		\end{align*}
		\begin{align*}
			\text{CRLB} = \frac{1}{nI(\mu)} = \frac{\sigma^2}{n} = \Var(\bar X). \quad \qedhere
		\end{align*}
	\end{block}
	\vspace{0.05cm}
	\pause
	\begin{alertblock}{Interpretation}
		\scriptsize
		Since $\bar X$ attains the CRLB exactly, it is the UMVUE for $\mu$: no unbiased estimator can have lower variance.
	\end{alertblock}
\end{frame}

% SLIDE 17: Exercise Level 4 (Statement)
\begin{frame}{In-Class Exercise --- Level 4: Challenging (1/2)}
	\vspace{-0.1cm}
	\begin{block}{Problem 6.1.9 --- Optimal Shrinkage Estimator (Statement)}
		With $\Var(\hat\theta)=v=5$ and $\theta=3$, numerically evaluate the optimal factor $c^*=\theta^2/(v+\theta^2)$ and compare $\text{MSE}(\hat\theta_{c^*})$ against $\text{MSE}(\hat\theta)=v$ (unbiased case).
	\end{block}
	\vspace{0.15cm}
	\pause
	\begin{alertblock}{Strategy}
		\scriptsize
		Substitute $v=5$, $\theta=3$ into $c^*$ and then into $\text{MSE}(c)=c^2v+(c-1)^2\theta^2$.
	\end{alertblock}
\end{frame}

% SLIDE 18: Exercise Level 4 (Solution)
\begin{frame}{In-Class Exercise --- Level 4: Challenging (2/2)}
	\vspace{-0.05cm}
	\scriptsize
	Substitute numerically: \pause
	\begin{block}{Calculation}
		\scriptsize
		\begin{align*}
			c^* = \frac{9}{5+9} = \frac{9}{14} \approx 0.643, \qquad \text{MSE}(c^*) \approx 2.066+1.148 = 3.214.
		\end{align*}
	\end{block}
	\vspace{0.05cm}
	\pause
	\begin{alertblock}{Interpretation}
		\scriptsize
		Compared to $\text{MSE}(c=1)=5$, the MSE decreases by roughly $35.7\%$ by accepting a controlled bias: a concrete example of why regularization (Ridge/Lasso) sacrifices unbiasedness for lower total error.
	\end{alertblock}
\end{frame}

% SLIDE 19: Closing and Chapter Perspective
\begin{frame}{Closing Section and Chapter Perspective}
	\vspace{-0.2cm}
	\begin{block}{Synthesis: Estimator Quality Criteria}
		\scriptsize
		Before learning to construct estimators (MoM, MLE), we needed a language to evaluate them. Bias, MSE, efficiency, and consistency are the four pillars we will use --- explicitly or implicitly --- throughout the rest of the chapter.
	\end{block}
	\vspace{0.05cm}
	\pause
	\begin{alertblock}{Key Lessons and Next Step}
		\begin{itemize}\itemsep=0.05cm
			\item \textbf{MSE = Variance + Bias$^2$:} a small bias can reduce total error.
			\item \textbf{Cramer-Rao:} $\Var(\hat\theta) \ge 1/(nI(\theta))$; attaining it defines the UMVUE.
			\item \textbf{Consistency:} if $\text{MSE}(\hat\theta_n)\to 0$, then $\hat\theta_n \xrightarrow{P} \theta$.
			\item \textbf{Next Section:} \textcolor{TecRojo}{Method of Moments (MoM)}.
		\end{itemize}
	\end{alertblock}
\end{frame}

\end{document}
